\documentclass[../main.tex]{subfiles}
\begin{document}
\begingroup
\emergencystretch=2em

\section{Experimental Setup}

This chapter evaluates whether the proposed Flip and ReFlip refinements improve quantization quality beyond standard round-to-nearest quantization and established post-training quantization baselines.

Two levels of evaluation are considered. The first is a standalone QKV experiment, where a selected grouped-query-attention block is isolated and the attention-score distortion is measured directly. This experiment makes the local effect of each refinement visible before running the full model. The second is a full-model evaluation, where quantized models are compared using perplexity and downstream task accuracy.

The standalone experiment follows the Transformer attention formulation \cite{vaswani2017attention} and focuses on the grouped-query-attention (GQA) structure used by modern decoder-only language models \cite{ainslie2023gqa}. It uses Llama-3-8B, layer 8, GQA group 3. Calibration activations are collected from C4 \cite{raffel2020t5} using 128 samples with sequence length 512, producing an activation matrix of size $65536 \times 4096$. The exported group contains four query heads sharing one key/value group. James-Stein shrinkage \cite{james1961estimation} is used to stabilize the activation mean estimate before quantization decisions are evaluated.

The full-model comparison covers Mistral-7B \cite{jiang2023mistral}, Llama-3/Llama-3.1 8B \cite{dubey2024llama3}, and Qwen2.5-7B \cite{qwen2024qwen25}. The results are taken from the 4-bit sheet of the reference workbook. Only the metrics used in this thesis are retained: WikiText-2 perplexity, C4 perplexity, ARC-Challenge, ARC-Easy, BoolQ, PIQA, RTE, and the average of the five accuracy tasks. WikiText-2 is a language modeling benchmark introduced with the WikiText corpus \cite{merity2016pointer}, while C4 originates from the T5 work \cite{raffel2020t5}. ARC-Challenge and ARC-Easy measure science-question reasoning \cite{clark2018arc}; BoolQ measures natural yes/no question answering \cite{clark2019boolq}; PIQA measures physical commonsense reasoning \cite{bisk2020piqa}; and RTE is included through the GLUE benchmark family \cite{wang2018glue}. The lm-evaluation harness provides a reproducible framework for running these tasks consistently across language models \cite{biderman2024reproducible}.

For WikiText-2 and C4, lower perplexity is better. For ARC-Challenge, ARC-Easy, BoolQ, PIQA, RTE, and their average, higher accuracy is better. The reported average is computed over the five selected downstream accuracy metrics:
\[
\mathrm{avg}=\frac{\mathrm{arc\mbox{-}c}+\mathrm{arc\mbox{-}e}+\mathrm{boolq}+\mathrm{piqa}+\mathrm{rte}}{5}.
\]
Full precision is shown as a reference ceiling rather than as a competitor in the quantized-method ranking. Therefore, in the full-model tables, bold values indicate the best quantized result for a metric within the same model, and underlined values indicate the second-best quantized result. Full-precision values are not highlighted.

The method names are normalized for clarity. The full-precision model is denoted as \textit{Full precision}. Standard round-to-nearest quantization is denoted as \textit{RTN}. RTN represents the direct rounding baseline commonly used in post-training quantization before more data-aware reconstruction is applied \cite{gholami2021survey,jacob2018quantization,nagel2020adaround}. GPTQ uses approximate second-order information to quantize pretrained Transformer weights \cite{frantar2022gptq}. AWQ is an activation-aware weight-only quantization method that protects salient channels using calibration statistics \cite{lin2023awq}. 

\section{Standalone QKV Evaluation}

The standalone QKV experiment compares three strategies under the same INT4 group-wise quantization setting. The nearest strategy performs direct round-to-nearest quantization. Flip then adjusts selected integer values by one quantization level to reduce local reconstruction error while avoiding activation outlier channels. ReFlip performs a second-stage correction that targets the attention-score error remaining after Flip.

\begin{table}[H]
\centering
\caption[Standalone QKV attention-score error]{Standalone QKV attention-score error on Llama-3-8B layer 8, GQA group 3.}
\label{tab:qkv-standalone-summary}
\setlength{\tabcolsep}{6pt}
\renewcommand{\arraystretch}{1.15}
\resizebox{\textwidth}{!}{%
\begin{tabular}{lrrrrr}
\hline
Method & Mean abs. error & Max abs. error & Mean rel. error (\%) & Reduction vs. RTN (\%) & Integer flips \\
\hline
RTN & 0.142806 & 0.253523 & 0.5330 & -- & -- \\
Flip & 0.070070 & 0.113184 & 0.2666 & 41.63 & 20,889 \\
ReFlip & \textbf{0.047939} & \textbf{0.068937} & \textbf{0.1934} & \textbf{56.63} & 64 \\
\hline
\end{tabular}%
}
\end{table}

Table \ref{tab:qkv-standalone-summary} shows that Flip substantially reduces the attention-score error compared with RTN, while ReFlip provides an additional reduction with very few extra integer changes. Flip modifies 20,889 weights across the selected Q and K projections, corresponding to approximately 0.7969\% of the involved weights. ReFlip applies only 64 additional integer changes, but still reduces the mean absolute attention-score error from 0.070070 to 0.047939.

\begin{table}[H]
\centering
\caption[Per-head QKV attention errors]{Per-head attention-score errors in the standalone QKV experiment.}
\label{tab:qkv-per-head}
\setlength{\tabcolsep}{5pt}
\renewcommand{\arraystretch}{1.15}
\resizebox{\textwidth}{!}{%
\begin{tabular}{lrrrrr}
\hline
Head & RTN error & Flip error & ReFlip error & RTN to Flip (\%) & RTN to ReFlip (\%) \\
\hline
0 & 0.092831 & 0.062511 & \textbf{0.034649} & 32.66 & 62.68 \\
1 & 0.253523 & 0.113184 & \textbf{0.068937} & 55.36 & 72.81 \\
2 & 0.192633 & 0.078560 & \textbf{0.063698} & 59.22 & 66.93 \\
3 & 0.032239 & 0.026023 & \textbf{0.024472} & 19.28 & 24.09 \\
\hline
\end{tabular}%
}
\end{table}

The per-head results in Table \ref{tab:qkv-per-head} indicate that ReFlip improves all four query heads in the selected GQA group. The magnitude of the improvement is not uniform across heads: head 1 obtains the largest RTN-to-ReFlip reduction, while head 3 shows a smaller but still positive reduction. This difference is expected because each query head has a different original score magnitude, error direction, and set of useful moderate dimensions for correction.

\begin{figure}[H]
\centering
\begin{subfigure}{0.48\textwidth}
    \centering
    \includegraphics[width=\linewidth]{../results/qkv_llama3_layer8_group3_full/method_summary_errors.png}
    \caption{Mean and maximum attention errors.}
\end{subfigure}
\hfill
\begin{subfigure}{0.48\textwidth}
    \centering
    \includegraphics[width=\linewidth]{../results/qkv_llama3_layer8_group3_full/manual_attention_error_by_head.png}
    \caption{Attention error by query head.}
\end{subfigure}
\caption[Standalone QKV error comparison]{Standalone QKV error comparison across RTN, Flip, and ReFlip.}
\label{fig:qkv-error-comparison}
\end{figure}

\begin{figure}[H]
\centering
\begin{subfigure}{0.48\textwidth}
    \centering
    \includegraphics[width=\linewidth]{../results/qkv_llama3_layer8_group3_full/error_waterfall_by_head.png}
    \caption{Error reduction from RTN to ReFlip.}
\end{subfigure}
\hfill
\begin{subfigure}{0.48\textwidth}
    \centering
    \includegraphics[width=\linewidth]{../results/qkv_llama3_layer8_group3_full/flip_cost_vs_error_reduction.png}
    \caption{Flip cost versus error reduction.}
\end{subfigure}
\caption[QKV correction behavior]{Per-head correction behavior and cost-benefit analysis.}
\label{fig:qkv-correction-behavior}
\end{figure}

The plots in Figures \ref{fig:qkv-error-comparison} and \ref{fig:qkv-correction-behavior} support the numerical tables. Flip provides a large first correction by changing a small fraction of weights, while ReFlip targets the remaining attention-score error more directly. The ReFlip stage therefore behaves as a fine correction step rather than a broad re-quantization pass.

\section{Full-Model Evaluation}

The full-model results evaluate whether the local improvements observed in the QKV setting are reflected in end-to-end model behavior. Four model families are included: Mistral-7B, Llama-3.1-8B, Qwen2.5-7B, and Llama-3-8B. For each model, the tables report 4-bit quantization results for RTN, AWQ, GPTQ, and their Flip/ReFlip variants when the corresponding workbook entries are available.

\begin{table}[H]
\centering
\caption[4-bit evaluation for Mistral-7B]{Full-model 4-bit evaluation for Mistral-7B. Bold marks the best quantized result; underline marks the second-best quantized result.}
\label{tab:full-model-mistral}
\setlength{\tabcolsep}{3pt}
\renewcommand{\arraystretch}{1.12}
\resizebox{\textwidth}{!}{%
\begin{tabular}{lrrrrrrrr}
\hline
Method & Wiki & C4 & arc-c & arc-e & boolq & piqa & rte & avg \\
\hline
Full precision & 4.845 & 7.535 & 0.509 & 0.802 & 0.837 & 0.800 & 0.686 & 0.7268 \\
RTN & 4.9917 & 7.8047 & 0.4761 & 0.7875 & \underline{0.8248} & 0.7992 & 0.6643 & 0.7104 \\
RTN + ReFlip & 4.9704 & \underline{7.7167} & 0.4804 & 0.7934 & 0.8202 & \underline{0.8014} & 0.6751 & 0.7141 \\
RTN + Flip & 4.9948 & 7.738 & 0.4744 & 0.7887 & \textbf{0.8275} & 0.7987 & 0.6606 & 0.7100 \\
AWQ & 4.975 & 7.727 & \underline{0.493} & 0.794 & 0.817 & 0.797 & 0.671 & 0.7144 \\
AWQ + ReFlip & \textbf{4.965} & \textbf{7.711} & \textbf{0.495} & \underline{0.795} & 0.818 & \textbf{0.802} & \textbf{0.690} & \textbf{0.7200} \\
AWQ + Flip & 4.975 & 7.726 & \underline{0.493} & 0.793 & 0.815 & 0.798 & \underline{0.678} & 0.7154 \\
GPTQ & 4.978 & 7.744 & 0.480 & \underline{0.795} & 0.820 & 0.798 & 0.675 & 0.7136 \\
GPTQ + ReFlip & \underline{4.967} & 7.736 & 0.483 & \textbf{0.796} & 0.821 & 0.800 & \underline{0.678} & \underline{0.7156} \\
GPTQ + Flip & 4.979 & 7.745 & 0.480 & \underline{0.795} & 0.819 & 0.798 & 0.675 & 0.7134 \\
\hline
\end{tabular}%
}
\end{table}

\begin{table}[H]
\centering
\caption[4-bit evaluation for Llama-3.1-8B]{Full-model 4-bit evaluation for Llama-3.1-8B. Bold marks the best quantized result; underline marks the second-best quantized result.}
\label{tab:full-model-llama31}
\setlength{\tabcolsep}{3pt}
\renewcommand{\arraystretch}{1.12}
\resizebox{\textwidth}{!}{%
\begin{tabular}{lrrrrrrrr}
\hline
Method & Wiki & C4 & arc-c & arc-e & boolq & piqa & rte & avg \\
\hline
Full precision & 5.522 & 8.409 & 0.522 & 0.822 & 0.830 & 0.791 & 0.711 & 0.7352 \\
RTN & 6.0248 & 9.4397 & 0.4898 & 0.8064 & \underline{0.7990} & 0.7835 & 0.6351 & 0.7147 \\
RTN + ReFlip & 5.9727 & 9.3692 & 0.5051 & \underline{0.8068} & 0.7928 & 0.7867 & 0.6531 & \textbf{0.7209} \\
RTN + Flip & 6.0815 & 9.6308 & 0.4872 & 0.8009 & \textbf{0.8075} & 0.7845 & 0.6315 & 0.7143 \\
AWQ & \underline{5.968} & 9.321 & \underline{0.511} & 0.806 & 0.794 & \textbf{0.798} & 0.635 & 0.7088 \\
AWQ + ReFlip & \textbf{5.949} & \textbf{9.294} & \textbf{0.512} & \textbf{0.808} & 0.796 & 0.796 & \textbf{0.671} & \underline{0.7166} \\
AWQ + Flip & \underline{5.968} & 9.316 & 0.507 & 0.803 & 0.790 & 0.796 & 0.632 & 0.7056 \\
GPTQ & 5.995 & 9.304 & 0.505 & 0.804 & 0.795 & 0.793 & 0.651 & 0.7096 \\
GPTQ + ReFlip & 5.989 & \underline{9.300} & 0.507 & 0.806 & 0.790 & \underline{0.797} & \underline{0.663} & 0.7126 \\
GPTQ + Flip & 5.992 & 9.304 & 0.499 & 0.803 & 0.792 & 0.793 & 0.659 & 0.7092 \\
\hline
\end{tabular}%
}
\end{table}

\begin{table}[H]
\centering
\caption[4-bit evaluation for Qwen2.5-7B]{Full-model 4-bit evaluation for Qwen2.5-7B. Bold marks the best quantized result; underline marks the second-best quantized result.}
\label{tab:full-model-qwen}
\setlength{\tabcolsep}{3pt}
\renewcommand{\arraystretch}{1.12}
\resizebox{\textwidth}{!}{%
\begin{tabular}{lrrrrrrrr}
\hline
Method & Wiki & C4 & arc-c & arc-e & boolq & piqa & rte & avg \\
\hline
Full precision & 6.0839 & 10.8166 & 0.477 & 0.804 & 0.846 & 0.787 & 0.815 & 0.7458 \\
RTN & 6.6915 & 12.2029 & 0.4506 & 0.7784 & 0.8302 & 0.7662 & 0.8016 & 0.7254 \\
RTN + ReFlip & 6.6366 & 11.8779 & 0.45162 & 0.77125 & 0.83264 & 0.76782 & 0.816018 & 0.7278696 \\
RTN + Flip & 6.6427 & 12.1076 & 0.44138 & 0.76832 & 0.82685 & 0.7651 & \textbf{0.823258} & 0.7249816 \\
AWQ & 6.4677 & \underline{11.4443} & \underline{0.469} & 0.786 & 0.842 & 0.780 & 0.805 & 0.7364 \\
AWQ + ReFlip & \textbf{6.438} & \textbf{11.363} & \textbf{0.473} & \textbf{0.792} & \textbf{0.847} & \underline{0.782} & \underline{0.819} & \textbf{0.7426} \\
AWQ + Flip & \underline{6.467} & 11.450 & 0.459 & \underline{0.787} & \underline{0.843} & \textbf{0.787} & 0.815 & \underline{0.7382} \\
GPTQ & 6.708 & 11.539 & 0.454 & 0.781 & 0.805 & 0.780 & 0.769 & 0.7178 \\
GPTQ + ReFlip & 6.702 & 11.534 & 0.459 & 0.779 & 0.814 & 0.779 & 0.787 & 0.7236 \\
GPTQ + Flip & 6.709 & 11.538 & 0.451 & 0.779 & 0.805 & 0.778 & 0.769 & 0.7164 \\
\hline
\end{tabular}%
}
\end{table}

\begin{table}[H]
\centering
\caption[4-bit evaluation for Llama-3-8B]{Full-model 4-bit evaluation for Llama-3-8B. Bold marks the best quantized result; underline marks the second-best quantized result.}
\label{tab:full-model-llama3}
\setlength{\tabcolsep}{3pt}
\renewcommand{\arraystretch}{1.12}
\resizebox{\textwidth}{!}{%
\begin{tabular}{lrrrrrrrr}
\hline
Method & Wiki & C4 & arc-c & arc-e & boolq & piqa & rte & avg \\
\hline
Full precision & 5.442 & 8.346 & 0.505 & 0.800 & 0.813 & 0.796 & 0.696 & 0.7220 \\
RTN & 6.0572 & 9.5846 & 0.4789 & 0.780 & 0.791 & 0.784 & 0.628 & 0.6924 \\
RTN + ReFlip & 6.0236 & 9.2523 & \underline{0.490} & 0.781 & 0.796 & 0.7818 & 0.6679 & 0.7033 \\
RTN + Flip & 6.0711 & 9.3279 & 0.4806 & 0.782 & 0.794 & 0.7802 & 0.6390 & 0.6952 \\
AWQ & 5.910 & 9.240 & 0.488 & 0.784 & 0.795 & \underline{0.791} & 0.646 & 0.7008 \\
AWQ + ReFlip & \textbf{5.885} & \textbf{9.199} & \textbf{0.493} & \underline{0.788} & \textbf{0.799} & 0.788 & 0.679 & \underline{0.7094} \\
AWQ + Flip & \underline{5.906} & \underline{9.234} & 0.489 & 0.784 & 0.796 & \textbf{0.792} & 0.617 & 0.6956 \\
GPTQ & 6.011 & 9.370 & 0.487 & 0.783 & \underline{0.798} & 0.788 & \underline{0.689} & 0.7090 \\
GPTQ + ReFlip & 6.013 & 9.354 & 0.489 & \textbf{0.790} & 0.795 & 0.788 & \textbf{0.693} & \textbf{0.7110} \\
GPTQ + Flip & 6.014 & 9.371 & 0.486 & 0.783 & 0.795 & 0.787 & 0.686 & 0.7074 \\
\hline
\end{tabular}%
}
\end{table}

\section{Discussion}

The standalone QKV experiment gives the clearest evidence for the mechanism targeted by this thesis. RTN introduces attention-score distortion because independent rounding errors in $Q$ and $K$ interact through the dot product. Flip reduces this error by improving local rounding decisions on non-outlier channels, and ReFlip further reduces the remaining error by explicitly correcting attention-score mismatch. The reduction from 0.142806 to 0.047939 in mean absolute attention error shows that the two-stage refinement is effective in the isolated setting. The cost profile is also important: Flip performs a broader local correction, while ReFlip adds only 64 integer flips and acts as a targeted residual correction step.

The full-model results are more nuanced because the downstream tasks depend on the entire network, not only on one QKV group. Nevertheless, several consistent patterns appear. On Mistral-7B, AWQ + ReFlip is the strongest quantized method for WikiText-2, C4, ARC-Challenge, PIQA, RTE, and the selected-task average. This suggests that the activation-aware scaling of AWQ and the attention-aware residual correction of ReFlip are complementary in this case. GPTQ + ReFlip obtains the best quantized ARC-Easy score, while RTN + Flip obtains the best quantized BoolQ score, showing that the best refinement can still vary by task.

On Llama-3.1-8B, AWQ + ReFlip gives the best quantized perplexity and the best quantized scores on ARC-Challenge, ARC-Easy, and RTE. RTN + ReFlip gives the highest quantized average, improving the RTN average from 0.7147 to 0.7209 while also reducing WikiText-2 and C4 perplexity. This indicates that ReFlip can be beneficial even when applied to the simplest RTN baseline, although the improvement is not uniformly distributed across every individual task.

On Qwen2.5-7B, AWQ + ReFlip is the strongest quantized method for WikiText-2, C4, ARC-Challenge, ARC-Easy, BoolQ, and the selected-task average. The RTN + ReFlip row also improves RTN perplexity and average task score, increasing the average from 0.7254 to 0.7279. RTN + Flip is particularly strong on RTE, reaching the best quantized RTE score in the table, but its average is slightly lower than RTN + ReFlip. This supports the interpretation that Flip and ReFlip correct different aspects of quantization error: Flip may strongly benefit selected tasks, while ReFlip tends to provide a more attention-aligned residual correction.

On Llama-3-8B, AWQ + ReFlip gives the best quantized WikiText-2 and C4 perplexity and also leads on ARC-Challenge and BoolQ. GPTQ + ReFlip gives the best quantized ARC-Easy, RTE, and average score. The improvements over each raw baseline are visible: RTN + ReFlip improves the RTN average from 0.6924 to 0.7033, AWQ + ReFlip improves AWQ from 0.7008 to 0.7094, and GPTQ + ReFlip improves GPTQ from 0.7090 to 0.7110. This is the clearest full-model evidence that ReFlip can improve multiple baseline quantizers rather than only a single quantization family.

Overall, ReFlip most consistently improves perplexity and average downstream performance when paired with AWQ or RTN, while GPTQ already provides a stronger reconstruction baseline and therefore leaves a smaller margin for further correction. Flip remains useful but is more task-dependent. This behavior is consistent with the design of the methods: Flip is a local rounding correction, whereas ReFlip explicitly targets the bilinear attention-score mismatch that arises from simultaneous Q and K quantization. The results should therefore be interpreted as evidence that attention-aware correction improves quantized model fidelity, not as a guarantee that every task score must improve after every refinement.

\endgroup
\end{document}
